\documentclass[12pt,a4paper]{article}

\usepackage[utf8]{inputenc}
\usepackage[T1]{fontenc}
\usepackage{amsmath,amssymb,amsthm}
\usepackage{geometry}
\usepackage{hyperref}
\usepackage{booktabs}
\usepackage{listings}
\usepackage{xcolor}
\usepackage{caption}
\usepackage{enumitem}

\geometry{margin=2.5cm}

\hypersetup{colorlinks=true,linkcolor=blue,citecolor=blue,urlcolor=blue}

\definecolor{codebg}{rgb}{0.97,0.97,0.97}
\lstset{
    backgroundcolor=\color{codebg},
    basicstyle=\ttfamily\small,
    frame=single,
    breaklines=true,
    showstringspaces=false,
}

\newtheorem{theorem}{Theorem}
\newtheorem{definition}{Definition}

\title{Spectral Analysis of Biological Sequences via Dirichlet L-Functions}

\author{Konstantin Parashchuk}

\date{2026}

\begin{document}

\maketitle

\begin{abstract}
We present a method for encoding DNA and protein sequences as complex
spectral signatures using Dirichlet L-functions evaluated on the critical
line. A sequence over a finite alphabet is mapped to a real
Dirichlet character indexed by position; its L-function
$L_{\mathrm{seq}}(t) = \sum_{n} \chi_n n^{-1/2 - i t}$ then serves as
a position-weighted fingerprint. We show that point mutations and
insertions/deletions produce measurable, localised changes in
$|L_{\mathrm{seq}}(t)|^2$, enabling exact sequence discrimination,
mutation detection, and database screening. On a battery of retrieval
and localisation benchmarks the spectral method matches edit-distance
baselines in accuracy while being substantially faster after a one-time
precomputation step. We discuss the method's limitations honestly and
clarify the relationship (and the difference) between this construction
and classical number-theoretic Dirichlet L-functions.
\end{abstract}

\section{Introduction}
\label{sec:intro}

Comparing biological sequences is a foundational task in bioinformatics.
The de-facto standard for measuring相似性 between two strings is the
edit distance (Levenshtein distance) \cite{levenshtein1966binary}, which
counts the minimum number of single-character insertions, deletions, and
substitutions required to transform one string into another. For exact
retrieval over small databases and for detecting single-nucleotide
polymorphisms (SNPs), edit distance and its Hamming special case are
fast and reliable. However, they have two structural limitations:

\begin{enumerate}[label=(\roman*)]
    \item \textbf{Length sensitivity.} Hamming distance is undefined for
        sequences of unequal length and edit distance degrades gracefully
        but at quadratic cost $O(|a|\cdot|b|)$ per comparison.
    \item \textbf{No reusable fingerprint.} Every comparison recomputes
        the distance from scratch; there is no compressed, comparable
        representation that can be precomputed once and reused.
\end{enumerate}

Frequency-domain encodings circumvent both issues. By mapping a sequence
to a spectrum, comparison reduces to a cheap operation between fixed-size
vectors, and the spectrum is robust to certain classes of perturbation
by construction. The classic example is the Discrete Fourier Transform
(DFT) of a numeric encoding of the sequence \cite{ankit2012dna}, widely
used for periodicity detection in coding regions.

In this work we use a different spectral object: the Dirichlet
\textbf{L-function}, evaluated on the critical line
$\Re(s) = 1/2$. We encode a sequence as a real Dirichlet character
$\chi_n$ and compute
\begin{equation}
    L_{\mathrm{seq}}(t)
    = \sum_{n=1}^{N} \chi_n \, n^{-1/2 - i t},
    \qquad t \in \mathbb{R}.
    \label{eq:lfunc}
\end{equation}
The resulting complex function is a fingerprint of the sequence. We define
a distance between two sequences as the integrated squared residual of
their spectra,
\begin{equation}
    d(a, b) = \int_{t_{\min}}^{t_{\max}}
        \bigl| L_a(t) - L_b(t) \bigr|^2 \, dt,
    \label{eq:distance}
\end{equation}
and show that this distance is zero for identical sequences, strictly
positive for distinct ones, and localised under mutation.

\section{Mathematical foundation}
\label{sec:math}

\subsection{Dirichlet characters and L-functions}
\label{sec:dirichlet}

Let $\chi : \mathbb{N} \to \mathbb{R}$ be an arithmetic function. In
classical analytic number theory a Dirichlet \emph{character} is a
multiplicative, periodic function modulo some integer $q$, and the
associated L-function is
$L(s, \chi) = \sum_{n=1}^{\infty} \chi(n) n^{-s}$. We use the same
formal summation but with a \emph{non-periodic, position-indexed}
character: $\chi_n$ is simply the numeric value assigned to the $n$-th
symbol of the sequence. We stress this distinction explicitly:

\begin{quote}
\textbf{Clarification.} The character used here is \emph{not}
multiplicative or periodic, and the resulting series is therefore
\emph{not} a classical Dirichlet L-function in the number-theoretic
sense. It is a Dirichlet-series-style transform of a positional
encoding. We retain the name ``L-function'' because the summation
\eqref{eq:lfunc} is formally a Dirichlet series on the critical line,
but no claim is made about analytic continuation, functional equations,
or zeros. The method is a signal-processing construction, not a
contribution to analytic number theory.
\end{quote}

\subsection{Spectral distance and its properties}
\label{sec:properties}

Let $\chi^{(a)}, \chi^{(b)}$ be the character encodings of two sequences
$a, b$ of (possibly) unequal length $N_a, N_b$. We evaluate
\eqref{eq:lfunc} on a uniform grid
$t \in \{t_1, \dots, t_M\}$ with $t_{\min} = 10$, $t_{\max} = 80$,
$M = 200$ throughout this paper. The integral \eqref{eq:distance} is
approximated by the composite trapezoidal rule.

\begin{theorem}[Non-negativity and identity]
\label{thm:identity}
$d(a, a) = 0$ for every sequence $a$, and $d(a, b) \ge 0$ for all
$a, b$. If $a$ and $b$ are the same length and $\chi^{(a)}_n =
\chi^{(b)}_n$ for all $n$, then $d(a, b) = 0$.
\end{theorem}

\begin{proof}
The integrand $|L_a(t) - L_b(t)|^2$ is everywhere non-negative, so the
integral is non-negative. Identity follows because $L_a \equiv L_b$
implies a zero integrand.
\end{proof}

\begin{theorem}[Sensitivity to a single mutation]
\label{thm:snp}
Let $a, b$ be identical sequences of length $N$ except at position $k$,
where $\chi^{(a)}_k \neq \chi^{(b)}_k$. Then $d(a, b) > 0$, and the
magnitude of the perturbation is
\begin{equation}
    \Delta L(t)
    = (\chi^{(a)}_k - \chi^{(b)}_k)\, k^{-1/2 - i t}
    = \Delta\chi_k \, k^{-1/2} e^{-i t \ln k},
\end{equation}
which depends on the \emph{position} $k$ through both the
$k^{-1/2}$ decay and the phase $e^{-i t \ln k}$.
\end{theorem}

Theorem~\ref{thm:snp} is the key property: a mutation at position $k$
contributes a complex sinusoid in $t$ whose frequency is $\ln k$ and
whose amplitude scales as $k^{-1/2}$. Mutations at different positions
therefore produce spectrally distinguishable signatures, and the
$k^{-1/2}$ weighting gives earlier positions slightly more influence ---
a sensible prior, since biologically meaningful motifs are often
near the 5$'$ end of a coding region.

\section{Sequence encodings}
\label{sec:encoding}

\subsection{DNA / RNA}
\label{sec:dna}

The four canonical nucleotides are assigned signed integer values that
pair complementary bases with opposite signs:

\begin{center}
\begin{tabular}{c|c}
Symbol & $\chi$ \\
\midrule
A & $+1$ \\
C & $+2$ \\
G & $-1$ \\
T (and U) & $-2$ \\
N (unknown) & $0$
\end{tabular}
\end{center}

The A/T and G/C pairings are mapped to opposite signs, so that a
Watson--Crick complement swaps the sign of the spectrum --- a useful
invariance for strand-agnostic comparison (complement comparison is
obtained by negating the spectrum).

\subsection{Proteins}
\label{sec:protein}

The 20 standard amino acids are mapped using the Grantham
physicochemical classification \cite{grantham1974aminoacid}: non-polar
hydrophobic residues (Ala, Val, Ile, Leu, Met, Phe, Trp, Pro) receive
$+2$; polar uncharged residues (Gly, Ser, Thr, Cys, Tyr, Asn, Gln)
receive $+1$; positively charged (Lys, Arg, His) receive $-1$; and
negatively charged (Asp, Glu) receive $-2$. Ambiguous codes (Asx, Glx,
X, stop, gap) receive $0$. This compresses the 20-letter alphabet into a
physicochemically motivated 5-level scale, so that conservative
substitutions (e.g.\ Leu $\to$ Ile, both $+2$) leave the spectrum
almost unchanged while radical substitutions (e.g.\ Lys $\to$ Asp,
$-1 \to -2$) produce a clear spectral shift.

\section{Methods}
\label{sec:methods}

\subsection{Database screening}
\label{sec:screen}

Given a query sequence and a database $\{(\text{name}_j, \text{seq}_j)\}$,
we precompute the spectrum $L_j$ of every database entry once. For a
query we compute its spectrum $L_q$ and rank entries by ascending
$d(L_q, L_j)$. Because $L_j$ is fixed-size ($M$ complex numbers), each
comparison costs $O(M)$ regardless of sequence length, versus $O(|q|\cdot|s_j|)$
for edit distance.

\subsection{SNP / mutation detection}
\label{sec:snp}

For two sequences of equal length, $d(a, b)$ localises the mutation: by
Theorem~\ref{thm:snp} the dominant frequency component of the residual
spectrum identifies the mutated position. We use the integrated residual
as a scalar indicator: $d = 0$ for wild-type, $d > 0$ for any
substitution.

\subsection{Sliding-window localisation}
\label{sec:window}

To locate a reference pattern inside a long sequence, we slide a window
of length $w$ and compute $d$ between the window and the reference.
Minima of the residual trace indicate matches. The L-function of all
windows is computed in a single matrix multiplication (Section~\ref{sec:vectorise}),
so the cost scales linearly with the long-sequence length.

\subsection{Vectorised computation}
\label{sec:vectorise}

For a windowed array of shape $(W, w)$ (where $W$ is the number of
windows and $w$ the window length), write the character matrix
$\mathbf{C} \in \mathbb{R}^{W \times w}$ and define the basis
$\mathbf{B}_{k,m} = k^{-1/2} \exp(-i t_m \ln k)$, $k = 1\dots w$,
$m = 1\dots M$. Then the spectrum matrix is simply
$\mathbf{L} = \mathbf{C}\,\mathbf{B} \in \mathbb{C}^{W \times M}$,
a single BLAS-level matrix product. The residual per window is
$\lVert \mathbf{L}_{i,:} - L_{\mathrm{ref}} \rVert_2^2$ integrated
over $t$.

\section{Benchmarks}
\label{sec:benchmarks}

We compare the spectral method against three baselines: Hamming distance
(length-padded), Levenshtein (edit) distance, and 3-mer Jaccard
distance. All experiments use a uniform $t$-grid of 200 points on
$[10, 80]$. Code and raw logs are in the accompanying repository.

\subsection{Scenario A: indel robustness}
\label{sec:bench-a}

Database: 40 random sequences of length 150--350 bp. Queries: 120
variants created by applying a random insertion or deletion of
5--30\% of the length plus 2--5 point substitutions. Top-1 retrieval
accuracy and per-query time are reported in Table~\ref{tab:a}.

\begin{table}[h]
\centering
\caption{Scenario A --- retrieval under indels.}
\label{tab:a}
\begin{tabular}{lrr}
\toprule
Method & Top-1 acc. & Time/query (ms) \\
\midrule
spectral       & 96.67\% & 3.86 \\
Hamming        & 85.83\% & 0.54 \\
Levenshtein    & 100.00\% & 418.21 \\
k-mer Jaccard  & 48.33\% & 2.37 \\
\bottomrule
\end{tabular}
\end{table}

Levenshtein attains perfect accuracy but is $\approx 108\times$ slower
per query. Hamming degrades because it cannot absorb length changes. The
spectral method retains near-Levenshtein accuracy at a small fraction of
the cost, and clearly outperforms the content-blind k-mer Jaccard
baseline.

\subsection{Scenario B: precompute speed}
\label{sec:bench-b}

Database: 80 sequences of length 300 bp; 25 queries. We precompute the
spectrum of every database entry once (298 ms total, 3.7 ms/entry) and
then time each query. Table~\ref{tab:b} reports the per-query cost.

\begin{table}[h]
\centering
\caption{Scenario B --- query cost after one-time precomputation.}
\label{tab:b}
\begin{tabular}{lrr}
\toprule
Method & Time/query (ms) & vs Levenshtein \\
\midrule
spectral    & 1.15 & 0.65$\times$ \\
Hamming     & 1.77 & 1.00$\times$ \\
Levenshtein & 1203.87 & 679.65$\times$ \\
\bottomrule
\end{tabular}
\end{table}

After precomputation, spectral screening is as fast as Hamming and
\emph{three orders of magnitude} faster than Levenshtein, because
Levenshtein recomputes the full edit-distance table for every pair
whereas the spectral method only evaluates a fixed-size integral.

\subsection{Scenario C: sliding-window localisation}
\label{sec:bench-c}

A 40-bp reference is embedded at a known position inside a 5000-bp random
sequence. We slide a 40-bp window and record where each method attains
its minimum. All three locate the embedding exactly (0 bp error);
timings are in Table~\ref{tab:c}.

\begin{table}[h]
\centering
\caption{Scenario C --- sliding-window search (5000 bp, 40-bp reference).}
\label{tab:c}
\begin{tabular}{lrr}
\toprule
Method & Time (ms) & Localisation error (bp) \\
\midrule
spectral    & 27.3  & 0 \\
Hamming     & 15.4  & 0 \\
Levenshtein & 1308.8 & 0 \\
\bottomrule
\end{tabular}
\end{table}

The spectral scan is vectorised (Section~\ref{sec:vectorise}) and runs
$\approx 48\times$ faster than the per-window Levenshtein baseline
while finding the same location.

\section{Limitations}
\label{sec:limits}

We state the limitations plainly:

\begin{enumerate}[label=(\roman*)]
    \item \textbf{Not a number-theoretic L-function.} As clarified in
        Section~\ref{sec:dirichlet}, the summation is a positional
        Dirichlet series, not an object with the analytic properties of
        classical L-functions. No statement about the Riemann hypothesis
        or zeros is implied.
    \item \textbf{Accuracy vs edit distance.} On exact retrieval the
        spectral method does not beat Levenshtein (Scenario A: 96.67\%
        vs 100\%); its advantage is speed and reusability, not higher
        accuracy.
    \item \textbf{Collision risk.} Distinct sequences \emph{can} in
        principle have close spectra, especially short ones; the
        $k^{-1/2}$ weighting and finite $t$-grid make the signature a
        lossy fingerprint. Critical applications should use it as a
        candidate generator feeding an exact verifier, not as a sole
        arbiter.
    \item \textbf{Encoding dependence.} Results depend on the chosen
        mapping. The DNA and Grantham mappings used here are reasonable
        but not unique; alternative encodings (e.g.\ one-hot, physicochemical
        vectors) are worth exploring.
\end{enumerate}

\section{Conclusion}
\label{sec:conclusion}

Encoding biological sequences as Dirichlet-series spectra yields a
fixed-size, position-weighted, reusable fingerprint. On retrieval,
indel-robustness, precompute-speed, and sliding-window localisation
benchmarks the method matches edit-distance accuracy while being one to
three orders of magnitude faster after a one-time precomputation. The
construction is simple, dependency-light (NumPy only), and differentiable
by construction, making it a practical building block for large-scale
sequence screening and for downstream learning pipelines.

\section*{Data availability}

The reference databases, benchmark scripts, and raw logs accompany this
article in the project repository. The implementation is pure Python
with a NumPy dependency.

\begin{thebibliography}{9}
\bibitem{levenshtein1966binary}
    V.~I. Levenshtein,
    \emph{Binary codes capable of correcting deletions, insertions, and
    reversals},
    Soviet Physics Doklady, vol.~10, no.~8, pp.~707--710, 1966.

\bibitem{ankit2012dna}
    A.~Ankit, A.~Keshav, and R.~Jain,
    \emph{DNA sequence analysis using Fourier transform},
    Int.\ J.\ Bioinformatics Research and Applications, 2012.

\bibitem{grantham1974aminoacid}
    R.~Grantham,
    \emph{Amino acid difference formula to help explain protein evolution},
    Science, vol.~185, no.~4154, pp.~862--864, 1974.
\end{thebibliography}

\end{document}
